% Rozdział 6 — szkic (2026-10-17), oparty na wynikach z rozdziału 5 (thesis/05-wyniki.tex)
% i na decyzjach zapisanych w CLAUDE.md. Bez odsyłaczy do innych rozdziałów — dodaj je przy redakcji.

\chapter{Podsumowanie i wnioski}
\label{ch:podsumowanie}

\section{Realizacja celu pracy}
\label{sec:podsumowanie-cel}

Celem pracy była analiza porównawcza relacyjnego systemu bazodanowego PostgreSQL wdrażanego
w~chmurze Microsoft Azure w~modelach IaaS i~PaaS, obejmująca wydajność, wykorzystanie zasobów
obliczeniowych oraz efektywność finansową. Cel zrealizowano w~trzech krokach. Po pierwsze,
zaprojektowano i~zautomatyzowano za pomocą Terraform dwa warianty środowisk: samodzielnie
skonfigurowany serwer PostgreSQL na maszynie wirtualnej z~dwoma tierami dysku (Standard SSD
i~Premium SSD) oraz usługę zarządzaną Azure Database for PostgreSQL -- Flexible Server
w~dwóch tierach obliczeniowych (Burstable i~General Purpose). Każde środowisko było tworzone
na czas sesji pomiarowej i~usuwane bezpośrednio po niej, a~konfiguracja PostgreSQL na
maszynie wirtualnej została wyrównana z~ustawieniami, które platforma stosuje w~usłudze
zarządzanej. Po drugie, przeprowadzono kampanię testów obciążeniowych narzędziem pgbench
(18 wdrożeń, 90 przebiegów w~stanie ustalonym), zbierając przepustowość, opóźnienia, IOPS
oraz wykorzystanie CPU i~pamięci. Po trzecie, połączono wyniki wydajności z~kosztem
utrzymania wyznaczonym z~cennika, uzyskując miarę efektywności finansowej w~postaci liczby
transakcji na dolara.

\section{Wnioski}
\label{sec:podsumowanie-wnioski}

\begin{enumerate}
  \item \textbf{Usługa zarządzana w~tierze General Purpose osiągnęła najwyższą wydajność.}
  Przepustowość PaaS General Purpose ($1358 \pm 19$~TPS) była o~45\% wyższa niż IaaS Premium
  SSD ($936 \pm 44$~TPS), a~99. percentyl opóźnienia niższy ($47$ wobec $69$~ms). Przewagę
  tę tłumaczą dwa czynniki ustalone w~eksperymencie wyjaśniającym: pamięć podręczna hosta,
  przez którą usługa zarządzana obsługuje większość odczytów poza limitem dysku, oraz klasa
  jednostki obliczeniowej -- 2~vCore bez mechanizmu kredytowego wobec maszyny
  \texttt{Standard\_B2s\_v2} z~poziomem bazowym 40\%. Samo włączenie pamięci podręcznej
  odczytu na maszynie wirtualnej usunęło wąskie gardło dyskowe, ale zwiększyło przepustowość
  tylko o~ok. 3,5\%, ponieważ ograniczeniem stał się poziom bazowy procesora.

  \item \textbf{Efektywność finansowa obu modeli w~konfiguracjach o~wysokiej wydajności jest
  porównywalna.} PaaS General Purpose i~IaaS Premium SSD dały odpowiednio $16{,}35$ i~$16{,}07$
  mln transakcji na dolara kosztu utrzymania, z~nakładającymi się przedziałami ufności.
  Wyższy koszt usługi zarządzanej (o~43\%) jest równoważony wyższą wydajnością, a~w~tej cenie
  usługa obejmuje kopie zapasowe z~odtwarzaniem do punktu w~czasie, aktualizacje
  i~monitorowanie, których koszt w~modelu IaaS ponosi administrator. Dla obciążeń
  transakcyjnych o~skali badanej w~pracy model PaaS nie jest zatem droższy w~przeliczeniu na
  wykonaną pracę.

  \item \textbf{W~modelu IaaS o~wydajności decyduje warstwa dyskowa.} Przy tej samej maszynie
  wirtualnej i~pojemności dysku Premium SSD P20 zapewnił $5{,}4$-krotnie wyższą przepustowość
  niż Standard SSD E20, a~serwer z~dyskiem E20 wykorzystywał procesor w~11\%. Dysk Standard SSD
  jest przy tym nieopłacalny pod intensywnym obciążeniem losowym: opłata za transakcje,
  naliczana do godzinowego limitu, prawie dwukrotnie przewyższa koszt jego pojemności, przez
  co konfiguracja z~dyskiem E20 jest droższa w~godzinie od konfiguracji z~P20 i~ma najniższą
  efektywność finansową spośród badanych ($2{,}46$~mln transakcji na dolara).

  \item \textbf{Tier Burstable usługi zarządzanej nadaje się do obciążeń okresowych, nie
  ciągłych.} PaaS Burstable jest najtańszy w~utrzymaniu, ale w~stanie utrzymywalnym osiąga
  $121$~TPS i~p99 opóźnienia powyżej sekundy, a~do tego stanu dochodzi dopiero po ponad
  trzech godzinach ciągłego obciążenia, gdy wyczerpią się kredyty CPU. Wcześniej przepustowość
  jest prawie dwukrotnie wyższa, co przy krótkich testach prowadzi do istotnego przeszacowania
  wydajności.

  \item \textbf{Mechanizmy kredytowe i~zmienność między wdrożeniami muszą być uwzględnione
  w~metodyce.} Pomiar wykonany bezpośrednio po wdrożeniu zawyżyłby wydajność IaaS Premium SSD
  o~ok. jedną trzecią. Zmienność wyników między kolejnymi wdrożeniami tej samej konfiguracji
  przewyższała w~części konfiguracji zmienność między powtórzeniami na jednym wdrożeniu, więc
  rzetelny przedział ufności wymaga wielu niezależnych wdrożeń -- co w~praktyce jest możliwe
  dzięki automatyzacji infrastruktury jako kodu i~niskiemu kosztowi środowisk efemerycznych
  (cała kampania z~pilotażem kosztowała ok. 20~USD).
\end{enumerate}

\section{Ograniczenia i zagrożenia trafności}
\label{sec:podsumowanie-ograniczenia}

\subsection*{Trafność wewnętrzna}

\begin{itemize}
  \item \textbf{Klasa maszyny wirtualnej.} Planowana maszyna \texttt{Standard\_D2s\_v5}
  (bez mechanizmu kredytowego) okazała się niedostępna dla subskrypcji Azure for Students,
  podobnie jak region Poland Central; użyto maszyny \texttt{Standard\_B2s\_v2} w~regionie
  Belgium Central. Porównanie IaaS z~PaaS General Purpose zestawia więc maszynę klasy
  burstable z~jednostką obliczeniową bez mechanizmu kredytowego, co -- jak pokazał
  eksperyment wyjaśniający -- odpowiada za część różnicy wydajności. Na głównych
  konfiguracjach IaaS kredyty CPU nie były wyczerpywane (procesor pracował poniżej poziomu
  bazowego), więc nie ograniczały wyników.
  \item \textbf{Pamięć podręczna hosta.} Dysk danych IaaS pracował bez pamięci podręcznej
  hosta, aby porównanie tierów dysku mierzyło warstwę dyskową, podczas gdy usługa zarządzana
  korzysta z~niej domyślnie. Wpływ tego wyboru oszacowano w~eksperymencie wyjaśniającym,
  wykonanym jednak tylko w~jednej ważnej sesji.
  \item \textbf{Mechanizmy kredytowe.} Stan pul kredytów kontrolowano metrykami Azure Monitor
  o~ograniczonej rozdzielczości (pula burst dysku: co 5~minut), a~dla usługi zarządzanej nie
  są publikowane metryki kredytów warstwy magazynu ani burstu I/O po stronie obliczeniowej.
  W~jednej sesji zaobserwowano reset kredytów dysku po stronie platformy; kryterium stanu
  ustalonego wykluczyło dotknięte nim przebiegi, ale procedura burn-inu potraktowała gwałtowny
  powrót puli jako stabilizację, a~nie sygnał do ponownego drenażu.
  \item \textbf{Parytet konfiguracji.} Na maszynę wirtualną przeniesiono wszystkie parametry
  PostgreSQL wpływające na wydajność, ale nie mechanizmy specyficzne dla platformy (ciągła
  archiwizacja WAL, rozszerzenia Azure, magazyn plików tymczasowych na dysku lokalnym). Ich
  narzut obciąża wyniki usługi zarządzanej i~jest zarazem ceną funkcji, które zapewnia.
  \item \textbf{Liczniki checkpointów.} Pierwotna kontrola wyrównania okna pomiaru z~cyklem
  checkpointu opierała się na liczniku \texttt{checkpoints\_timed}, publikowanym z~opóźnieniem,
  gdy zapis nie nadąża z~harmonogramem; zastąpiono ją liczbą zakończonych checkpointów, co
  potwierdziło jeden cykl w~89 z~90 okien.
  \item \textbf{Awarie infrastruktury pomiarowej.} Kampania była sterowana z~lokalnego
  komputera; chwilowy brak miejsca na jego dysku uniemożliwił usunięcie jednego środowiska,
  a~restart usługi SSH w~trakcie aktualizacji pakietów i~chwilowy brak mocy obliczeniowej
  w~regionie przerwały kilka prób wdrożenia. Próby nieudane powtórzono, a~przebiegi sesji
  przerwanej pominięto w~analizie (analiza wrażliwości pokazała, że ich dołączenie nie zmienia
  wniosków).
\end{itemize}

\subsection*{Trafność zewnętrzna}

\begin{itemize}
  \item \textbf{Jedno obciążenie.} Wyniki dotyczą obciążenia TPC-B-podobnego narzędzia pgbench
  (losowe, krótkie transakcje zapisujące) dla bazy 14,6~GiB i~25~klientów. Obciążenia
  zdominowane przez odczyty, analityczne lub z~bazą mieszczącą się w~pamięci mogą dać inną
  kolejność konfiguracji -- w~szczególności inną relację kosztu dysku Standard SSD, którego
  opłata transakcyjna zależy od liczby operacji.
  \item \textbf{Jeden region, jeden okres, jedna subskrypcja.} Pomiary wykonano w~regionie
  Belgium Central w~październiku 2026~r.; wydajność współdzielonej infrastruktury może
  zależeć od regionu i~obciążenia platformy w~czasie. Kolejność i~pora sesji były losowane
  i~rotowane, a~zmienność między wdrożeniami ujęto w~przedziałach ufności, ale nie obejmuje
  ona innych regionów ani dłuższych okresów.
  \item \textbf{Rozmiar magazynu.} Pojemność 512~GiB przyjęto we wszystkich konfiguracjach dla
  zachowania jednakowej warstwy dyskowej (przy pojemności 128~GB tiery Standard i~Premium
  mają identyczny limit 500~IOPS, więc porównanie warstwy dyskowej nie miałoby czego mierzyć). Dla PaaS Burstable magazyn stanowi ponad 80\% kosztu, więc
  przy mniejszej bazie relacja kosztów byłaby inna.
\end{itemize}

\subsection*{Model kosztów}

\begin{itemize}
  \item Koszty wyznaczono z~cennika katalogowego pay-as-you-go z~17~października 2026~r.,
  bez rabatów za rezerwację, planów oszczędnościowych i~bezpłatnych limitów; ceny chmury się
  zmieniają, dlatego migawkę cennika zapisano w~repozytorium.
  \item Rzeczywisty rachunek subskrypcji nie mógł posłużyć do weryfikacji kosztu usługi
  zarządzanej, ponieważ dane Cost Management nie zawierały żadnej pozycji Azure Database for
  PostgreSQL. Dla zasobów IaaS rachunek był zgodny z~cennikiem, w~tym z~limitem naliczania
  transakcji dysku Standard SSD.
  \item Model obejmuje wyłącznie koszt zasobów strony bazy danych; nie uwzględnia kosztu pracy
  administracyjnej, transferu danych ani wysokiej dostępności, które w~praktyce różnicują
  modele IaaS i~PaaS na korzyść usługi zarządzanej.
\end{itemize}

\section{Kierunki dalszych prac}
\label{sec:podsumowanie-dalsze}

\begin{itemize}
  \item Powtórzenie porównania z~maszyną wirtualną bez mechanizmu kredytowego (np. seria
  Dsv5), aby oddzielić wpływ modelu wdrożenia od klasy jednostki obliczeniowej.
  \item Rozszerzenie o~inne profile obciążenia (tylko odczyt, obciążenie mieszane, większa skala
  bazy i~liczba klientów) oraz o~dyski Premium SSD v2 z~niezależnie konfigurowanymi IOPS.
  \item Uwzględnienie konfiguracji wysokiej dostępności (replikacja, strefy dostępności)
  i~kosztu ich utrzymania w~obu modelach.
  \item Model kosztów uwzględniający rezerwacje i~plany oszczędnościowe oraz koszt pracy
  administracyjnej (całkowity koszt posiadania).
  \item Uzupełnienie procedury burn-inu o~wykrywanie gwałtownego odnowienia puli kredytów
  i~ponowny drenaż w~takim przypadku.
\end{itemize}
